\section{Introduction}\label{sec:introduction}

A \emph{hypersymplectic structure} on an oriented four-manifold is a triple
of closed real two-forms whose pointwise span is three-dimensional and
positive definite for the wedge product. Such a triple determines a conformal structure,
but its forms need not be parallel. After a constant normalization of
their intersection matrix, the question is whether one can deform the
forms to a hyperk\"ahler triple while retaining their three cohomology
classes.

\begin{theorem}\label{thm:main}
Let $X$ be a closed connected oriented smooth four-manifold, and let
$\omega=(\omega_1,\omega_2,\omega_3)$ be a smooth hypersymplectic structure
normalized by
\begin{equation}\label{eq:intro-normalization}
 \int_X\omega_i\wedge\omega_j=\delta_{ij}.
\end{equation}
There is a smooth family of hypersymplectic structures $\omega(t)$,
$0\le t\le1$, such that
\[
 \omega(0)=\omega,\qquad [\omega_i(t)]=[\omega_i]
 \quad (i=1,2,3),
\]
and a positive volume form $\mu_1$ with
\begin{equation}\label{eq:intro-endpoint}
 \omega_i(1)\wedge\omega_j(1)=2\delta_{ij}\mu_1.
\end{equation}
The endpoint forms are parallel and self-dual for a hyperk\"ahler metric.
\end{theorem}

For an arbitrary positive volume form $\mu$, write
$\omega_i\wedge\omega_j=2\mathcal Q_{ij}\mu$. There is a unique choice
of $\mu$ for which $\det\mathcal Q=1$, and
\eqref{eq:intro-endpoint} is exactly $\mathcal Q=I_3$ in this convention.
The statement is unaffected by a constant orthogonal change of the
ordered triple.

The normalization also specifies the conclusion for any initial positive
closed triple. Let $G_{ij}=\int_X\omega_i\wedge\omega_j$; this matrix is
positive definite. Choose a constant invertible real matrix $L$ with
$LGL^{\mathsf T}=I_3$ and apply \Cref{thm:main} to $L\omega$.
Applying $L^{-1}$ to the resulting path preserves positivity and each
original class. Its endpoint satisfies
\[
 \omega_i(1)\wedge\omega_j(1)=2G_{ij}\mu_1,
\]
where $\mu_1$ is the volume form of the normalized endpoint metric.
Thus the original basis consists of parallel self-dual forms for a
hyperk\"ahler metric, with its original Gram matrix $G$.

\begin{corollary}[Individual K\"ahler realizations]\label{cor:individual-kahler}
Let $X$ be a closed connected oriented smooth four-manifold with a smooth
positive triple $\omega$ of closed real two-forms. For every fixed
$c\in\mathbb R^3\setminus\{0\}$, the form $\alpha_c=\sum_i c_i\omega_i$
is a K\"ahler form for a parallel integrable complex structure $J_c$ of
some hyperk\"ahler metric $g_c$ on $X$:
\[
 \alpha_c(v,w)=g_c(J_cv,w).
\]
The compatible structure is chosen separately for each $c$; no simultaneous
transport of the triple is asserted.
\end{corollary}
\begin{proof}
Use the fixed-class path for the original basis described above and put
$\alpha_t=\sum_i c_i\omega_i(t)$. Positivity of the pointwise wedge-product
matrix gives $\alpha_t\wedge\alpha_t>0$. Thus $\alpha_t$ is symplectic,
and its class is fixed because each $[\omega_i(t)]$ is fixed.

For the endpoint metric $g$, set $r_c=(c^{\mathsf T}Gc)^{1/2}>0$.
The endpoint identities show that $\alpha_1/r_c$ is parallel and self-dual
with $g$-norm $\sqrt2$. It therefore defines a parallel integrable complex
structure $J^{\mathrm{end}}_c$ with
$\alpha_1/r_c=g(J^{\mathrm{end}}_c\,\cdot,\cdot)$.
Hence $\alpha_1$ is its K\"ahler form for the hyperk\"ahler metric $r_cg$.

Since $X$ is closed, Moser's theorem gives an isotopy $\phi_{c,t}$ with
$\phi_{c,0}=\mathrm{id}$ and $\phi_{c,t}^*\alpha_t=\alpha_c$.
Pulling back $r_cg$ and $J^{\mathrm{end}}_c$ by $\phi_{c,1}$ gives the
claimed hyperk\"ahler metric and compatible integrable structure on $X$.
Fine and Yao note this Moser implication for a constituent in
\cite[Section~1]{FineYao2017}.
\end{proof}

\begin{corollary}[Smooth type]\label{cor:smooth-type}
Let $X$ be a closed connected smooth four-manifold admitting a smooth
positive triple of closed real two-forms. Then $X$ is diffeomorphic to
the standard K3 manifold or to $T^4$.
\end{corollary}
\begin{proof}
Give $X$ the orientation induced by the triple. The normalization above
and \Cref{thm:main} give a global hyperk\"ahler triple on $X$. Choosing
one of its complex structures makes $X$ a compact K\"ahler surface with
a nowhere-vanishing holomorphic two-form. The compact hyperk\"ahler
surface classification \cite[Section~2, p.~162]{Boyer1988} therefore
identifies $X$ itself as a complex two-torus or a K3 surface. A complex
two-torus is smoothly $T^4$, while every complex K3 surface is
diffeomorphic to a nonsingular quartic surface in $\mathbf P^3$, the
standard K3 manifold \cite[p.~58]{Kodaira1970}.
\end{proof}

\subsection{History and significance}

Donaldson's study of closed two-forms on four-manifolds links their
pointwise wedge-product geometry to nonlinear elliptic equations
\cite{Donaldson2006}. His Question~3 in Section~5.3 asks whether a
compact oriented four-manifold carrying a positive closed triple
admits a hyperk\"ahler structure. In the same section he proposes a
prescribed-volume continuity argument in the simply connected case,
conditional on a priori estimates and an auxiliary involution. The
involution keeps the evolving form in a fixed cohomology class
\cite[Section~5.3]{Donaldson2006}.
Fine and Yao formulate the normalized
deformation problem, retaining all three cohomology classes, as
Conjecture~1.1 of \cite{FineYao2017}.
\Cref{thm:main} resolves this deformation conjecture positively.
Keeping the classes connects the given triple to a metric with its
prescribed periods, beyond existence of some hyperk\"ahler metric on
the same manifold.

The main analytic approach has been the hypersymplectic flow.
Fine and Yao obtained this flow by reducing the $G_2$-Laplacian flow
on the product with a three-torus. They proved that a finite-time
solution extends while the associated seven-dimensional scalar
curvature, equivalently the torsion, remains bounded
\cite[Theorem~1.3]{FineYao2017}. Their later work gives an integral
torsion extension criterion, global existence near pointwise
orthogonality, and subsequential convergence of global solutions under
additional geometric bounds
\cite[Theorems~4.7, 4.10, 4.15 and~4.16]{FineYao2020}.

Several restricted classes admit complete convergence results.
Huang, Wang and Yao proved convergence, after pullback by
diffeomorphisms, for simple-type hypersymplectic structures on
the standard four-torus, a diagonal one-variable class
\cite[Theorems~3.5 and~3.6]{HuangWangYao2018}.
Picard and Suan treated $G_2$ flows associated with K\"ahler
Calabi--Yau data; their complex-dimension-two case supplies a
hypersymplectic convergence theorem
\cite[Theorem~1.3]{PicardSuan2024}.
Fine, He and Yao extended the simple-type convergence result, modulo
diffeomorphisms, to
$T^3$-invariant triples on $T^4$ in symmetric normal form
\cite[Theorem~1.5]{FineHeYao2024}. They subsequently proved
a linear cohomologous isotopy to a hyperk\"ahler triple when an
effective smooth circle action preserves the initial triple
\cite[Theorem~1.2]{FineHeYao2025}.
Related reductions and coflow constructions appear in work of Petcu,
Yao and Zhou, and Karigiannis, Picard and Suan
\cite{Petcu2026,YaoZhou2026,KarigiannisPicardSuan2026}.
These results retain their stated symmetry or integrability hypotheses.
The path constructed here starts from an arbitrary positive closed
triple; its existence is a separate question from convergence or
uniqueness for the hypersymplectic flow.

The topology already places strong restrictions on such a triple.
Any two of its forms determine an almost-complex structure whose
canonical line is trivial, and the third form tames that structure.
Here a real two-form $\beta$ \emph{tames} $J$ if
$\beta(v,Jv)>0$ for every nonzero vector $v$; it is
\emph{compatible} if it is also $J$-invariant. A closed tamer is
symplectic. Bauer's bound for symplectic four-manifolds with torsion
first Chern class \cite[Corollary~1.2]{Bauer2008}, together with the
three positive period classes, gives $b^+=3$, where $b^+$ is the
positive index of the intersection form. The resulting free
intersection lattice is the K3 lattice or three hyperbolic planes.
These facts supply the lattice used below; they do not presuppose a
smooth identification of the initial manifold with a K3 surface or torus.

For controlling closed positive forms, Sullivan's structure currents
provide the cone-duality framework \cite{Sullivan1976}.
Harvey and Lawson developed the positive-current characterization of
K\"ahler geometry \cite{HarveyLawson1983}, and Li and Zhang compared
the taming and compatible cones in almost-complex geometry
assuming the compatible cone is nonempty \cite{LiZhang2009}.
A separate article proves that every tamed smooth
almost-complex structure on a closed connected real four-manifold admits a compatible
symplectic form \cite[Theorem~1.1]{OpenAITaming2026}. Its current
estimates and density-splitting argument are also used here, with the
needed proofs included in \Cref{sec:current-energy,sec:density-splitting}.
In particular, smooth density-weight closedness already appears in
the proof in Section~6 of that article. Its compatible class is
unrestricted. The additional class criterion below converts closed
normalized density thresholds into curves, using Preiss's
rectifiability theorem and Rivi\`ere--Tian's regularity theorem for
integral cycles \cite{Preiss1987,RiviereTian2009}. This supplies the
prescribed taming class needed along the deformation.

\subsection{The main ideas}

The proof organizes the deformation around a pair of forms. Write an
ordered triple as $(A,C,B)$, and write products of forms for wedge
products. The positive plane spanned by $A,C$
determines an almost-complex structure $J$, with sign chosen so that
$B$ tames it. A pair deformation gives the desired triple deformation
provided the original third class continues to contain a tamer.
The argument first proves a criterion for this availability, then
constructs a fibration and an auxiliary equal-square pair, and finally
uses two classical K\"ahler steps to reach the endpoint.

\paragraph{Currents and a prescribed class.}
Separation of the cone of positive forms gives a positive current
obstructing a desired closed form. Quadratic mass estimates bound its
mass in a ball of radius $r$ by $Kr^2$ for a fixed constant $K$.
Resolvent regularization also
controls how its complex-line directions vary on that scale.
Together these estimates prove closedness after weighting by smooth
functions of the two-dimensional density. The zero-density part can
then be handled by the intersection form, while normalized thresholds
of the positive-density part give integral pseudoholomorphic cycles.
Compatibility is proved before the integral-cycle regularity theorem
is applied, so that its local symplectic hypothesis is available.

The output is \Cref{thm:tamed-compatible,thm:taming-cone}.
For the second theorem, let $V$ be the intersection-orthogonal
complement of $[A],[C]$. It has Lorentz signature because $b^+=3$.
A positive-square class $H\in V$ is a taming class if it is in the
same positive component as a known positive-square taming class $U\in V$ and
pairs positively with every irreducible closed $J$-curve. The
criterion will keep $[B]$ available as the pair changes.

\paragraph{Marked curves and a torus fibration.}
The initial triple supplies an entire coefficient sphere of
almost-complex structures: for $v\in S^2$, the forms
$a\cdot\omega$ with $a\cdot v=0$ constitute its anti-invariant plane
(a form $\xi$ is anti-invariant when $\xi(Ju,Jw)=-\xi(u,w)$),
and $v\cdot\omega$ tames the
corresponding structure. Compatible forms can be selected smoothly
over this sphere by a constant-rank elliptic correction. After projecting
their classes to the original positive period plane and normalizing,
one obtains a degree-one map to its unit sphere.
Li proposed using such winding sphere families, parametrized wall crossing,
and Taubes's curve correspondence to construct elliptic fibrations
\cite[Section~7.4.1]{Li2010SCY}. The marked-family construction below
carries out this approach for the coefficient sphere of the initial triple.
Li--Liu's families wall-crossing and index formulas
\cite{LiLiu2001}, together with Taubes's canonical nonvanishing and
large-parameter curve extraction
\cite{Taubes1994,Taubes1996,Taubes1999}, then give curves through every
point in a chosen nonzero integral square-zero class.
To impose passage through a chosen point, the proof requires one
component of the Seiberg--Witten spinor to vanish there. As the parameter moves
over the sphere, the line containing that component also varies and
has degree of absolute value one. Its contribution must be retained
in the marked wall-crossing calculation, even though the underlying
Spin-c structure is fixed throughout the family.

The integral class is chosen to prevent any such curve configuration
from splitting. Eichler's criterion supplies the lattice normal form
\cite[Proposition~3.3(i)]{GritsenkoHulekSankaran2009}; a further
finite-coset argument works for the arbitrary real positive period
plane of the original triple. Pseudoholomorphic adjunction and
positivity of intersections then reduce the possible fibers to tori
and rational curves with one singularity
\cite{McDuffSalamon2012,Wendl2020}.
To obtain ordinary nodes, we prove transversality using exact
perturbations of $A,C$ with $B$ fixed. The first-jet method is related
to that of Oh and Zhu \cite{OhZhu2008}, but the allowable closed-form
variations require their own cokernel calculation. Relative nodal
replacement and local deformation theory now produce the fibration
of \Cref{thm:pencil}: a proper map from $X$ to a closed oriented
surface, with torus regular fibers and rational singular fibers having
one ordinary node.

\paragraph{An auxiliary volume equation.}
The fibration supplies classes with small positive fiber area.
Near a nodal fiber we use the Gibbons--Hawking construction
\cite{GibbonsHawking1978} and the periodic Ooguri--Vafa model
\cite{OoguriVafa1996}. Gross and Wilson developed the corresponding
mathematical collapsing and semi-flat gluing methods for elliptic
K3 surfaces \cite{GrossWilson2000}. Here the nodal holomorphic model
is constructed independently, since a global complex structure is
not yet available. On fixed annuli the model and regular data agree
up to exponentially small error. Matching the fiber area and the
period on the surface swept out by the monodromy-invariant fiber
cycle ensures that the discrepancy has a primitive.

After exact preparation of the pair, these local forms glue to an
approximate auxiliary form of fiber area $\epsilon$.
The correction in \Cref{thm:auxiliary-solution} is solved on exact
two-forms. For any Riemannian metric on the oriented closed four-manifold, let
$h^+$ denote the self-dual part of an exact two-form $h$. The exact-form
identity used in Donaldson's elliptic formulation
\cite[Section~2, proof of Proposition~1; Section~4, Lemma~1]{Donaldson2006}
\[
 \|h\|_{L^2}^2=2\|h^+\|_{L^2}^2\qquad(h\text{ exact})
\]
controls the inverse while the fibers collapse. The remaining finite
Sobolev estimates lose only powers of $\epsilon^{-1}$, which are
absorbed by the exponentially small gluing error. This produces a
closed form $D$ satisfying $AD=CD=0$ and $D^2=A^2$.
The class inequality in \Cref{prop:auxiliary-class} holds for all
later curves at one fixed choice of $\epsilon$; it transfers
positivity from $[D]$ to the original class $[B]$.

\paragraph{Two K\"ahler steps in the original classes.}
The closed equal-square pair $A,D$ defines an integrable complex
structure with holomorphic two-form $A\pm iD$, choosing the sign
so that $C$ tames it. The compact-surface
criteria of Buchdahl and Lamari, and the numerical K\"ahler-cone
criterion, put $[C]$ in the K\"ahler cone
\cite{Buchdahl1999,Lamari1999,DemaillyPaun2004}.
Yau's volume theorem \cite{Yau1978} gives $C_1\in[C]$ with
$C_1^2=A^2$. Interpolate from $C$ to $C_1$; the uniform class
inequality and the taming criterion keep $[B]$ available throughout.
The pair $A,C_1$ is now itself an integrable equal-square pair.
A second K\"ahler step in $[B]$ gives $B_1$, completing the
hyperk\"ahler triple. Smooth selection of the third forms, with both
endpoints prescribed, gives the actual path $(A,C_t,B_t)$ shown in
\Cref{fig:endpoint-path}. Every component retains its original class.

\subsection{Organization and conventions}

\Cref{sec:definite-pairs} establishes the preliminary geometry and
topology. The positive-current results occupy
\Cref{sec:current-energy,sec:density-splitting}.
\Cref{sec:families,sec:lattice-chamber,sec:pencil} construct the
torus fibration, and
\Cref{sec:local-models,sec:regular-preparation,sec:auxiliary-volume}
construct the auxiliary equal-square pair. The proof of
\Cref{thm:main} is completed in \Cref{sec:endpoint}.

We write products of two-forms for wedge products and products of
degree-two cohomology classes for their intersection pairing.
Curve classes are identified with their Poincar\'e duals in cohomology.
Integral lattice statements use
\[
 \Lambda=H^2(X;\Z)/\operatorname{torsion}.
\]
The original positive period plane is
$P=\operatorname{span}\{[\omega_1],[\omega_2],[\omega_3]\}$.
No rationality of this plane is assumed. All deformations of forms are
smooth in the manifold variable; parameter regularity is stated where
it is used. Constants in local estimates may depend on fixed smooth
background data and on the indicated finite differentiation order.
